% Discussion
\newpage
\section{Discussion}
\label{sec:discussion}

\subsection{Comparison with Prior Multi-Agent Systems}
\label{subsec:comparison}

Table~\ref{tab:comparison} contrasts SPOQ with representative
multi-agent systems and autonomous coding tools.

\begin{table}[htbp]
\centering
\caption{Comparison of Multi-Agent and Autonomous Coding Approaches}
\label{tab:comparison}
\footnotesize
\begin{tabular}{lccccccc}
\toprule
\textbf{Feature} & \textbf{SPOQ} & \textbf{ChatDev} & \textbf{MetaGPT} & \textbf{AutoGPT} & \textbf{Devin} & \textbf{Aider} \\
\midrule
Execution model & Wave-parallel & Sequential & Sequential & Recursive & Single-agent & Single-agent \\
Explicit dependencies & DAG & Implicit & Implicit & Priority queue & None & None \\
Planning validation & 10 metrics & None & Informal & None & None & None \\
Code validation & 10 metrics & Test-based & Review & Self-eval & Test-based & Test-based \\
Human integration & HaaA & Observer & Observer & Minimal & Minimal & Collaborative \\
Agent specialization & 3-tier & Role-based & Role-based & General & General & General \\
Parallelism potential & High & Low & Low & Low & None & None \\
\bottomrule
\end{tabular}
\end{table}

\paragraph{Execution Model.}
ChatDev and MetaGPT use sequential role-playing where agents take turns.
MAD uses synchronous debate rounds. SPOQ achieves true parallelism through
wave-based dispatch, with agents in the same wave executing concurrently.

\paragraph{Dependency Management.}
Prior systems embed dependencies implicitly in role sequences (programmer
after designer) or artifact flows. SPOQ makes dependencies explicit in a
DAG, enabling static analysis, critical path computation, and parallelism
optimization.

\paragraph{Validation Rigor.}
ChatDev relies on test execution; MetaGPT uses informal review.
SPOQ applies structured metrics at both planning and code stages,
with quantified thresholds that prevent low-quality work from proceeding.

\paragraph{Human Role.}
In ChatDev and MetaGPT, humans observe outputs but do not participate
in the orchestration loop. SPOQ's HaaA model integrates human judgment
at key decision points while preserving automation benefits.

\subsection{Limitations}
\label{subsec:limitations}

SPOQ has several limitations that future work should address:

\paragraph{Upfront Planning Investment.}
SPOQ requires detailed epic specifications before execution begins.
For exploratory or rapidly-changing requirements, this upfront cost
may be prohibitive. Projects with unclear scope may benefit from
more adaptive approaches.

\paragraph{Dependency on Human Quality.}
The HaaA model's effectiveness depends on the human specialist's
skill in task decomposition and validation. A human who approves
poorly-structured epics will see cascading quality issues.

\paragraph{Empirical Scale.}
Our controlled benchmark suite covers four full-stack tasks across two
provider families (Claude and a locally hosted Qwen3.6-35B-A3B). The
accompanying deployment study spans 17 repositories, 1{,}822 completed
tasks, and 13{,}866 executed tests, with one third-party adopter. While
this breadth strengthens generalizability claims relative to a single
case study, the four-task benchmark suite is still modest, the Qwen
replication is single-shot rather than mean$\pm$sd, and most field
deployments share a single primary author. Independent replication by
other teams, larger benchmark suites with repeated stochastic runs,
and broader controlled studies with matched baselines are needed to
validate speedup and quality claims rigorously.

\paragraph{Reference Implementation Coupling.}
Our current implementation is coupled to Claude Code and Anthropic's
model family, representing a specific instantiation of the SPOQ methodology
rather than its only possible form. However, the core methodology, including wave
computation, validation scoring, dependency resolution, and journal
tracking, is expressed in platform-agnostic YAML and Markdown formats.
Porting SPOQ to alternative platforms (Gemini CLI, OpenAI Assistants API,
open-source model deployments) would require remapping the capability
tiers to the target provider's model offerings and adapting the
orchestration interface. The fundamental algorithms and quality gates
remain unchanged across implementations.

\paragraph{No Cross-Epic Learning.}
Each epic starts fresh. SPOQ does not currently transfer lessons
from prior epics (successful patterns, common failure modes) to
new planning sessions.

\paragraph{Metric Reliability.}
The 20 validation metrics rely on LLM-based assessment without
inter-rater reliability testing. Automated quality scores may
exhibit biases, inconsistency across runs, or susceptibility
to gaming. Human validation studies comparing LLM scores to
expert assessments are needed.

\paragraph{Cost Approximations.}
Our cost analysis uses estimated token counts that vary
significantly by task complexity, codebase size, and context
requirements. Real-world costs may differ substantially. The
analysis does not include infrastructure, human oversight time,
or rework loop costs comprehensively.

\subsection{Future Work}
\label{subsec:future}

A platform-level development addresses SPOQ's orchestration overhead
and defines the near-term trajectory of the methodology.

\paragraph{Native Platform Support: Claude Code Agent Teams.}
SPOQ's current orchestration relies on manually-coordinated Claude Code
instances. Claude Code's experimental Agent Teams feature provides native
primitives that align directly with SPOQ's dispatch model: \texttt{TeamCreate}
establishes a coordinated session with a lead and specialized teammates,
\texttt{TaskCreate}/\texttt{TaskUpdate} manage a shared task list with
dependency tracking, and \texttt{SendMessage} enables direct inter-agent
communication. Teams support plan approval workflows where teammates
operate in read-only mode until the lead validates their approach, mirroring
SPOQ's validation gates.

The mapping is natural: SPOQ's wave computation produces task assignments
with explicit dependencies; Agent Teams' task list enforces those
dependencies through blocked/unblocked state transitions. SPOQ's three-tier
hierarchy (Opus workers, Haiku investigators, Sonnet reviewers) maps to
Agent Teams' model selection per teammate. The delegate mode, which
restricts the lead to coordination-only tools, formalizes SPOQ's
orchestrator role separation.

Integrating SPOQ with Agent Teams would shift the methodology from a
set of conventions enforced by skill prompts to a native orchestration
protocol. Wave dispatch becomes \texttt{TaskCreate} with dependency
declarations. Validation gates become \texttt{TaskCompleted} hooks that
invoke Sonnet reviewers before marking tasks done. The journal system
feeds naturally into the shared task list's completion records.
This creates a pipeline where the human declares a goal, the
orchestrator generates a wave-structured task list, and Agent Teams
execute it with built-in coordination, messaging, and quality enforcement.

\subsection{Broader Implications}
\label{subsec:implications}

SPOQ represents a step toward \textit{AI-native software engineering}:
development processes designed around AI capabilities rather than
retrofitting AI into human-centric workflows.

Key implications include:

\begin{itemize}
    \item \textbf{Redefined roles:} Engineers shift from writing code
    to validating agent outputs, decomposing problems, and making
    architectural decisions.

    \item \textbf{Quality as constraint:} Explicit quality gates force
    upfront investment in planning, potentially improving overall
    software quality.

    \item \textbf{Scalable expertise:} A single human specialist can
    leverage multiple agents, scaling their expertise across more
    projects than traditional pair programming allows.
\end{itemize}

\paragraph{The Director Model.}
At scale, SPOQ enables a single engineer to direct a digital workforce:
6 Claude Code instances (1 planning copilot + 5 execution overseers)
coordinating 50--100 concurrent agents across independent epics, achieving
daily throughput of 75--150 tasks with output previously requiring
8--10 engineers.

\paragraph{Applicability.}
SPOQ's structured approach introduces upfront overhead that pays off for
projects with 5+ parallelizable tasks spanning 4+ hours of estimated work.
Below this threshold, orchestration overhead typically exceeds execution
time savings. Appendix~\ref{subsec:usage-guidance-detail} provides a
detailed decision framework with applicability guidance by project type.
